\uuid{lhCW}
\exo7id{6069}
\titre{exo7 6069}
\auteur{queffelec}
\organisation{exo7}
\datecreate{2011-10-16}
\isIndication{false}
\isCorrection{false}
\chapitre{Continuité, uniforme continuité}
\sousChapitre{Continuité, uniforme continuité}
\cpge{mpsi-mp2i.limites-continuite.b, pcsi.limites-continuite.b, ptsi.limites-continuite.b, mpsi-mp2i.fonctions-valeurs-reelles-complexes.a, pcsi.techniques-calcul-fonctions.a, ptsi.techniques-calcul-fonctions.a, bcpst1.limites-continuite.c, bcpst1.fonctions-reelles-usuelles.a, mp-mpi.topologie-espaces-normes.e, pc.espaces-vectoriels-normes.e, psi.espaces-vectoriels-normes.e, pt.fonctions-deux-trois-variables.b}
\module{Topologie}
\niveau{L3}
\difficulte{}

\contenu{
\texte{
Soit $f$ une application de $\R$ dans $\R$ et $\omega$ sa fonction
oscillation définie pour $x_0\in \R$ et $\delta>0$ par 
$$\omega(x_0,\delta)=\sup_{\{|x_0-y|=\delta,|x_0-z|=\delta\}}|f(y)-f(z)|.$$
}
\begin{enumerate}
    \item \question{Remarquer que $f$ est continue en $x_0$ si et seulement si
$$\omega(x_0)=\inf_{\delta>0}\omega(x_0,\delta)=0.$$}
    \item \question{Montrer que pour tout $\varepsilon >0$, $O_\varepsilon=\{x\ ;\
\omega(x)<\varepsilon\}$ est un ouvert.

En déduire que $C(f)$, l'ensemble des points de continuité de $f$, est un
$G_\delta$.}
\end{enumerate}
}
